\section{Self-Attention: de la alineación entre dos secuencias a las relaciones dentro de una secuencia}

La encoder--decoder attention tradicional suele hacer que la target-side query lea las source-side keys y values; la self-attention, en cambio, hace que query, key y value provengan del mismo conjunto de representaciones. Este cambio parece consistir solo en pasar el mismo tensor como los tres argumentos de la función, pero su significado estadístico cambia: el modelo ya no se limita a alinear secuencias distintas, sino que empieza a aprender un grafo de relaciones dentro de una misma muestra.

\lecturefigure{slide-03-self-attention.jpg}{La self-attention define source y target como la misma secuencia y aprende intra-sequence relations.}{Video oficial de Stanford Online, 00:01:42--00:02:39.}

Sea la representación de entrada $X\in\mathbb{R}^{m\times d_{\mathrm{model}}}$; la scaled dot-product attention se escribe como
\[
Q=XW_Q,\qquad K=XW_K,\qquad V=XW_V,
\]
\[
A=\operatorname{softmax}\!\left(\frac{QK^{\top}}{\sqrt{d_k}}\right),
\qquad Z=AV.
\]

\begin{itemize}
  \item $m$: longitud de la secuencia, es decir, el número de tokens que participan en el modelado de relaciones;
  \item $W_Q,W_K,W_V$: matrices aprendibles que proyectan la entrada en query, key y value;
  \item $d_k$: dimensión de cada key/query; $\sqrt{d_k}$ sirve para controlar la escala del dot product;
  \item $A_{ij}$: peso normalizado con el que el token $i$ lee información del token $j$;
  \item $Z_i$: nueva representación del token $i$, obtenida al agregar todos los values.
\end{itemize}

\teachervoice{El ponente usa «blue ball» para ilustrar la semántica de la self-attention: el modelo debería aprender a asociar el adjetivo blue con el sustantivo ball al que modifica. Este ejemplo no afirma que cierta head corresponda necesariamente a un árbol sintáctico; subraya que los attention weights pueden expresar hipótesis sobre las relaciones internas de la entrada.}

\begin{warningbox}{Los attention weights no son una explicación que se obtenga automáticamente}
$A_{ij}$ es el peso con el que se mezcla la información en una pasada hacia adelante. Puede indicar dónde está leyendo el modelo, pero por sí solo no demuestra una relación lingüística, una relación causal ni el motivo de una decisión. Para interpretarlo hay que considerar también el contenido de los values, las capas posteriores, la ablation o la intervention.
\end{warningbox}

\subsection{Multi-Head Attention: aprendizaje en paralelo de distintos subespacios de relaciones}

Un solo attention map solo puede comparar query y key en un único espacio de proyección. La multi-head attention proyecta la representación en varios subespacios de menor dimensión, lo que permite que distintas heads formen distintas medidas de similitud; después concatena los resultados y los vuelve a mezclar. Al leer la figura, conviene seguir primero el flujo de datos «división--attention independiente--concatenación», en lugar de entender el número de heads como una simple copia del modelo.

\lecturefigure{slide-04-multi-head-attention.jpg}{Multi-head attention: varios subespacios de relaciones se calculan de forma independiente y luego se concatenan.}{Video oficial de Stanford Online, 00:02:40--00:03:27.}

\[
\operatorname{head}_r
=\operatorname{Att}\!\left(QW_r^Q,KW_r^K,VW_r^V\right),
\]
\[
\operatorname{MHA}(Q,K,V)
=\operatorname{Concat}(\operatorname{head}_1,\ldots,\operatorname{head}_H)W^O.
\]

\begin{itemize}
  \item $H$: número de attention heads;
  \item $W_r^Q,W_r^K,W_r^V$: proyecciones independientes de la head $r$;
  \item $W^O$: vuelve a mezclar las salidas de todas las heads en la model dimension;
  \item cada head puede aprender relaciones distintas, pero ningún mecanismo garantiza una división del trabajo perfecta y nombrable por una persona.
\end{itemize}

\teachervoice{La intuición de Gomez es que una head puede inclinarse hacia la relación adjetivo--sustantivo y otra puede aprender otros patrones. Esto debe leerse como «la capacidad permite repartir las relaciones», no como «después del entrenamiento cada head se convierte automáticamente en un concepto humano claro».}

\begin{knowledgebox}{Por qué se divide entre $\sqrt{d_k}$}
Si las dimensiones de query y key son aproximadamente independientes y de varianza constante, la varianza del dot product crece con $d_k$. Los logits demasiado grandes hacen que la softmax se sature prematuramente y que los gradientes se reduzcan; el factor de escala devuelve los valores a un intervalo más estable.
\end{knowledgebox}

\subsection{El Transformer no es un único truco}

La self-attention y la multi-head attention son los títulos fáciles de recordar, pero Gomez enfatiza que el sistema original también dependía de un conjunto de decisiones de entrenamiento que en su momento no eran evidentes: LayerNorm, residual path, learning-rate warmup, la inicialización y los valores predeterminados de ingeniería de Tensor2Tensor. Tras la publicación del artículo, estos elementos se volvieron conocimiento común, lo que puede llevar a los lectores posteriores a creer que estos componentes, por naturaleza, debían combinarse así.

\teachervoice{El equipo realizó ablations repetidamente y encontró que retirar el warmup o la LayerNorm, o modificar configuraciones aparentemente menores, deterioraba de verdad la optimización. La conclusión de ingeniería es que el éxito de una arquitectura suele provenir de un conjunto de mecanismos que funcionan en conjunto, y que no se puede reducir toda la contribución a una sola fórmula de attention.}

\teachervoice{Gomez recuerda que el trabajo original tomó forma rápidamente en unos tres meses y se terminó a toda prisa antes de la fecha límite; en ese momento no previó su impacto. La lectura más valiosa de esta historia no consiste en fabricar un «momento de genialidad», sino en reconocer que las herramientas maduras, la iteración entre varias personas, la retroalimentación experimental y la oportunidad del momento contribuyeron juntas a un nuevo paradigma.}

\subsection{Resumen de la sección}

La self-attention traslada el modelado de relaciones de entre dos secuencias al interior de una misma entrada; la multi-head attention aporta varios subespacios de relaciones. La capacidad real de entrenar el Transformer depende además de un conjunto de técnicas de optimización y de conexiones residuales, y no puede explicarse solo con la expresión $QK^{\top}V$.
